\section{Задание}

Решить систему $AX=B$ методами Якоби и Зейделя с точностью
\[
\varepsilon = 10^{-4}.
\]

\medskip
\noindent\textbf{Вариант 47:}

\[
A=
\begin{pmatrix}
5{,}02 & 0{,}13 & 0{,}21 & 0{,}08\\
0{,}11 & 6{,}04 & 0{,}14 & 0{,}17\\
0{,}09 & 0{,}12 & 7{,}01 & 0{,}10\\
0{,}13 & 0{,}08 & 0{,}16 & 8{,}03
\end{pmatrix},
\qquad
B=
\begin{pmatrix}
10{,}54\\
-5{,}40\\
21{,}09\\
0{,}66
\end{pmatrix}.
\]

Требуется найти столбец
\[
X=
\begin{pmatrix}
x_1\\
x_2\\
x_3\\
x_4
\end{pmatrix}.
\]

\section{Цель работы}

Научиться применять методы Якоби и Зейделя для решения систем линейных алгебраических уравнений.

\section{Ход работы}

Рассмотрим систему
\[
AX=B.
\]

В развёрнутом виде она имеет вид:
\[
\begin{cases}
5{,}02x_1+0{,}13x_2+0{,}21x_3+0{,}08x_4=10{,}54,\\
0{,}11x_1+6{,}04x_2+0{,}14x_3+0{,}17x_4=-5{,}40,\\
0{,}09x_1+0{,}12x_2+7{,}01x_3+0{,}10x_4=21{,}09,\\
0{,}13x_1+0{,}08x_2+0{,}16x_3+8{,}03x_4=0{,}66.
\end{cases}
\]

\subsection{Проверка условия сходимости}

Проверим диагональное преобладание матрицы $A$.

Для первой строки:
\[
|a_{11}|=5{,}02,
\qquad
|a_{12}|+|a_{13}|+|a_{14}|=0{,}13+0{,}21+0{,}08=0{,}42.
\]

Для второй строки:
\[
|a_{22}|=6{,}04,
\qquad
|a_{21}|+|a_{23}|+|a_{24}|=0{,}11+0{,}14+0{,}17=0{,}42.
\]

Для третьей строки:
\[
|a_{33}|=7{,}01,
\qquad
|a_{31}|+|a_{32}|+|a_{34}|=0{,}09+0{,}12+0{,}10=0{,}31.
\]

Для четвёртой строки:
\[
|a_{44}|=8{,}03,
\qquad
|a_{41}|+|a_{42}|+|a_{43}|=0{,}13+0{,}08+0{,}16=0{,}37.
\]

Так как
\[
|a_{ii}|>\sum_{j\ne i}|a_{ij}|,
\]
то матрица $A$ имеет строгое диагональное преобладание. Следовательно, методы Якоби и Зейделя сходятся.

\subsection{Приведение системы к итерационному виду}

Выразим каждую неизвестную через остальные:

\[
\begin{cases}
x_1=\dfrac{10{,}54-0{,}13x_2-0{,}21x_3-0{,}08x_4}{5{,}02},\\[8pt]
x_2=\dfrac{-5{,}40-0{,}11x_1-0{,}14x_3-0{,}17x_4}{6{,}04},\\[8pt]
x_3=\dfrac{21{,}09-0{,}09x_1-0{,}12x_2-0{,}10x_4}{7{,}01},\\[8pt]
x_4=\dfrac{0{,}66-0{,}13x_1-0{,}08x_2-0{,}16x_3}{8{,}03}.
\end{cases}
\]

После вычисления коэффициентов получаем:
\[
\begin{cases}
x_1=2{,}099602-0{,}025896x_2-0{,}041833x_3-0{,}015936x_4,\\
x_2=-0{,}894040-0{,}018212x_1-0{,}023179x_3-0{,}028146x_4,\\
x_3=3{,}008559-0{,}012839x_1-0{,}017118x_2-0{,}014265x_4,\\
x_4=0{,}082192-0{,}016189x_1-0{,}009963x_2-0{,}019925x_3.
\end{cases}
\]

Таким образом, система приведена к виду
\[
X=C+\alpha X,
\]
где
\[
C=
\begin{pmatrix}
2{,}099602\\
-0{,}894040\\
3{,}008559\\
0{,}082192
\end{pmatrix},
\]

\[
\alpha=
\begin{pmatrix}
0 & -0{,}025896 & -0{,}041833 & -0{,}015936\\
-0{,}018212 & 0 & -0{,}023179 & -0{,}028146\\
-0{,}012839 & -0{,}017118 & 0 & -0{,}014265\\
-0{,}016189 & -0{,}009963 & -0{,}019925 & 0
\end{pmatrix}.
\]

В качестве начального приближения возьмём
\[
X^{(0)}=
\begin{pmatrix}
0\\
0\\
0\\
0
\end{pmatrix}.
\]

Вычислим норму матрицы $\alpha$:
\[
\|\alpha\|_{\infty}
=
\max_i \sum_{j=1}^{4}|\alpha_{ij}|.
\]

Получаем:
\[
\|\alpha\|_{\infty}=0{,}083665<1.
\]

Следовательно, итерационный процесс сходится.

Для остановки итерационного процесса используем апостериорную оценку:
\[
\frac{\|\alpha\|}{1-\|\alpha\|}
\cdot
\|X^{(k)}-X^{(k-1)}\|
\leq
\varepsilon.
\]

Отсюда
\[
\|X^{(k)}-X^{(k-1)}\|
\leq
\frac{1-\|\alpha\|}{\|\alpha\|}\varepsilon.
\]

Подставим значения:
\[
\varepsilon_1=
\frac{1-0{,}083665}{0{,}083665}\cdot 10^{-4}
\approx 0{,}001095.
\]

Значит, итерации продолжаем до выполнения условия
\[
\|X^{(k)}-X^{(k-1)}\|_{\infty}\leq 0{,}001095.
\]

\subsection{Решение методом Якоби}

Метод Якоби имеет вид:
\[
x_i^{(k)}
=
\frac{1}{a_{ii}}
\left(
b_i-\sum_{j\ne i}a_{ij}x_j^{(k-1)}
\right).
\]

Для данной системы получаем расчётные формулы:
\[
\begin{cases}
x_1^{(k)}=2{,}099602-0{,}025896x_2^{(k-1)}-0{,}041833x_3^{(k-1)}-0{,}015936x_4^{(k-1)},\\
x_2^{(k)}=-0{,}894040-0{,}018212x_1^{(k-1)}-0{,}023179x_3^{(k-1)}-0{,}028146x_4^{(k-1)},\\
x_3^{(k)}=3{,}008559-0{,}012839x_1^{(k-1)}-0{,}017118x_2^{(k-1)}-0{,}014265x_4^{(k-1)},\\
x_4^{(k)}=0{,}082192-0{,}016189x_1^{(k-1)}-0{,}009963x_2^{(k-1)}-0{,}019925x_3^{(k-1)}.
\end{cases}
\]

Выполним итерации.

\begin{table}[h!]
\centering
\caption{Последовательность приближений метода Якоби}
\begin{tabular}{|c|c|c|c|c|c|}
\hline
$k$ & $x_1^{(k)}$ & $x_2^{(k)}$ & $x_3^{(k)}$ & $x_4^{(k)}$ & $\|X^{(k)}-X^{(k-1)}\|_{\infty}$\\
\hline
0 & 0{,}000000 & 0{,}000000 & 0{,}000000 & 0{,}000000 & --\\
\hline
1 & 2{,}099602 & -0{,}894040 & 3{,}008559 & 0{,}082192 & 3{,}008559\\
\hline
2 & 1{,}995588 & -1{,}004326 & 2{,}995735 & -0{,}002839 & 0{,}110286\\
\hline
3 & 2{,}000336 & -0{,}999741 & 3{,}000171 & 0{,}000200 & 0{,}004748\\
\hline
4 & 1{,}999983 & -1{,}000016 & 2{,}999988 & -0{,}000011 & 0{,}000353\\
\hline
\end{tabular}
\end{table}

На четвёртой итерации:
\[
\|X^{(4)}-X^{(3)}\|_{\infty}=0{,}000353.
\]

Так как
\[
0{,}000353<0{,}001095,
\]
то условие остановки выполнено.

Следовательно, приближённое решение, найденное методом Якоби:
\[
X_{\text{Якоби}}
\approx
\begin{pmatrix}
1{,}999983\\
-1{,}000016\\
2{,}999988\\
-0{,}000011
\end{pmatrix}.
\]

Оценим погрешность по апостериорной формуле:
\[
\Delta X
\leq
\frac{\|\alpha\|_{\infty}}{1-\|\alpha\|_{\infty}}
\cdot
\|X^{(4)}-X^{(3)}\|_{\infty}.
\]

\[
\Delta X
\leq
\frac{0{,}083665}{1-0{,}083665}\cdot0{,}000353
\approx
0{,}000032.
\]

Полученная оценка меньше заданной точности:
\[
0{,}000032<10^{-4}.
\]

\subsection{Решение методом Зейделя}

Метод Зейделя отличается от метода Якоби тем, что при вычислении очередной компоненты используются уже найденные значения текущей итерации.

Расчётные формулы имеют вид:
\[
\begin{cases}
x_1^{(k)}=2{,}099602-0{,}025896x_2^{(k-1)}-0{,}041833x_3^{(k-1)}-0{,}015936x_4^{(k-1)},\\
x_2^{(k)}=-0{,}894040-0{,}018212x_1^{(k)}-0{,}023179x_3^{(k-1)}-0{,}028146x_4^{(k-1)},\\
x_3^{(k)}=3{,}008559-0{,}012839x_1^{(k)}-0{,}017118x_2^{(k)}-0{,}014265x_4^{(k-1)},\\
x_4^{(k)}=0{,}082192-0{,}016189x_1^{(k)}-0{,}009963x_2^{(k)}-0{,}019925x_3^{(k)}.
\end{cases}
\]

Выполним итерации.

\begin{table}[h!]
\centering
\caption{Последовательность приближений метода Зейделя}
\begin{tabular}{|c|c|c|c|c|c|}
\hline
$k$ & $x_1^{(k)}$ & $x_2^{(k)}$ & $x_3^{(k)}$ & $x_4^{(k)}$ & $\|X^{(k)}-X^{(k-1)}\|_{\infty}$\\
\hline
0 & 0{,}000000 & 0{,}000000 & 0{,}000000 & 0{,}000000 & --\\
\hline
1 & 2{,}099602 & -0{,}932278 & 2{,}997562 & -0{,}002239 & 2{,}997562\\
\hline
2 & 1{,}998384 & -0{,}999851 & 3{,}000050 & 0{,}000024 & 0{,}101218\\
\hline
3 & 1{,}999994 & -1{,}000002 & 3{,}000000 & 0{,}000000 & 0{,}001610\\
\hline
4 & 2{,}000000 & -1{,}000000 & 3{,}000000 & -0{,}000000 & 0{,}000006\\
\hline
\end{tabular}
\end{table}

На четвёртой итерации:
\[
\|X^{(4)}-X^{(3)}\|_{\infty}=0{,}000006.
\]

Так как
\[
0{,}000006<0{,}001095,
\]
то условие остановки выполнено.

Следовательно, приближённое решение, найденное методом Зейделя:
\[
X_{\text{Зейделя}}
\approx
\begin{pmatrix}
2{,}000000\\
-1{,}000000\\
3{,}000000\\
0{,}000000
\end{pmatrix}.
\]

Оценим погрешность.

Для метода Зейделя матрицу $\alpha$ представим в виде
\[
\alpha=\alpha_1+\alpha_2,
\]
где $\alpha_1$ --- нижняя треугольная часть, а $\alpha_2$ --- верхняя треугольная часть.

\[
\alpha_1=
\begin{pmatrix}
0 & 0 & 0 & 0\\
-0{,}018212 & 0 & 0 & 0\\
-0{,}012839 & -0{,}017118 & 0 & 0\\
-0{,}016189 & -0{,}009963 & -0{,}019925 & 0
\end{pmatrix},
\]

\[
\alpha_2=
\begin{pmatrix}
0 & -0{,}025896 & -0{,}041833 & -0{,}015936\\
0 & 0 & -0{,}023179 & -0{,}028146\\
0 & 0 & 0 & -0{,}014265\\
0 & 0 & 0 & 0
\end{pmatrix}.
\]

Нормы этих матриц равны:
\[
\|\alpha_1\|_{\infty}=0{,}046077,
\qquad
\|\alpha_2\|_{\infty}=0{,}083665.
\]

Оценка погрешности:
\[
\Delta X
\leq
\frac{\|\alpha_2\|_{\infty}}{1-\|\alpha\|_{\infty}}
\cdot
\|X^{(4)}-X^{(3)}\|_{\infty}.
\]

Подставим значения:
\[
\Delta X
\leq
\frac{0{,}083665}{1-0{,}083665}\cdot0{,}000006
\approx
0{,}000001.
\]

Полученная оценка меньше заданной точности:
\[
0{,}000001<10^{-4}.
\]

\subsection{Проверка решения}

Полученное решение:
\[
X\approx
\begin{pmatrix}
2\\
-1\\
3\\
0
\end{pmatrix}.
\]

Подставим его в систему:

\[
AX=
\begin{pmatrix}
5{,}02 & 0{,}13 & 0{,}21 & 0{,}08\\
0{,}11 & 6{,}04 & 0{,}14 & 0{,}17\\
0{,}09 & 0{,}12 & 7{,}01 & 0{,}10\\
0{,}13 & 0{,}08 & 0{,}16 & 8{,}03
\end{pmatrix}
\begin{pmatrix}
2\\
-1\\
3\\
0
\end{pmatrix}.
\]

Выполним умножение:

\[
\begin{aligned}
5{,}02\cdot2+0{,}13\cdot(-1)+0{,}21\cdot3+0{,}08\cdot0
&=10{,}04-0{,}13+0{,}63=10{,}54,\\
0{,}11\cdot2+6{,}04\cdot(-1)+0{,}14\cdot3+0{,}17\cdot0
&=0{,}22-6{,}04+0{,}42=-5{,}40,\\
0{,}09\cdot2+0{,}12\cdot(-1)+7{,}01\cdot3+0{,}10\cdot0
&=0{,}18-0{,}12+21{,}03=21{,}09,\\
0{,}13\cdot2+0{,}08\cdot(-1)+0{,}16\cdot3+8{,}03\cdot0
&=0{,}26-0{,}08+0{,}48=0{,}66.
\end{aligned}
\]

Следовательно,
\[
AX=
\begin{pmatrix}
10{,}54\\
-5{,}40\\
21{,}09\\
0{,}66
\end{pmatrix}
=B.
\]

Значит, найденное решение корректно.

\section{Вывод}

В ходе лабораторной работы была решена система линейных алгебраических уравнений
\[
AX=B
\]
методами Якоби и Зейделя с точностью
\[
\varepsilon=10^{-4}.
\]

Для варианта 47 было получено решение:
\[
X=
\begin{pmatrix}
2\\
-1\\
3\\
0
\end{pmatrix}.
\]

Метод Якоби дал приближённое решение:
\[
X_{\text{Якоби}}
\approx
\begin{pmatrix}
1{,}999983\\
-1{,}000016\\
2{,}999988\\
-0{,}000011
\end{pmatrix}.
\]

Метод Зейделя дал приближённое решение:
\[
X_{\text{Зейделя}}
\approx
\begin{pmatrix}
2{,}000000\\
-1{,}000000\\
3{,}000000\\
0{,}000000
\end{pmatrix}.
\]

Оба метода сошлись, так как матрица системы имеет строгое диагональное преобладание. Метод Зейделя в данной задаче оказался точнее на той же итерации и быстрее приблизился к точному решению.
